\documentclass{article}
\usepackage[T1]{fontenc}
\usepackage{lmodern}
\usepackage{graphicx}
\usepackage{amsmath,amssymb}
\usepackage[a4paper,margin=2cm]{geometry}
\usepackage[absolute,overlay]{textpos}
\setlength{\TPHorizModule}{1mm}
\setlength{\TPVertModule}{1mm}
\usepackage[indonesian]{babel}
\usepackage{hyperref}
\setlength{\emergencystretch}{2em}
% unit: Halaman 26

\begin{document}

% === Halaman 26 (python/native) ===

\{ Dekripsi RC6-w/r/b

Masukan: Cipherteks disimpan di dalam empat register berukuran

w bit: A, B, C dan D r adalah jumlah putaran

S[0, ... , 2r + 3] adalah kunci internal, panjangnya w bit

Luaran: Plainteks disimpan di dalam A, B, C, D

\}

Algoritma:

C C – S[2r + 3]

A A – S[2r + 2]

for i r downto 1 do

(A, B, C, D) (D, A, B, C)

u (D*(2D + 1)) <<< lg w

t (B*(2B + 1)) <<< lg w

C ((C – S[2i + 1]) >>> t) u


      A  ((A – S[2i]) >>> u) t

end

D D – S[1]

B B – S[0]

26

\end{document}
